\subsection{Feature Ablation Study}

To assess the contribution of engineered features and identify critical predictors, we conducted a systematic feature ablation study. Table \ref{tab:feature_ablation} presents model performance when removing specific feature groups.

\begin{table}[htbp]
\centering
\caption{Feature Ablation Study: Impact of Removing Feature Groups}
\label{tab:feature_ablation}
\begin{tabular}{lcccc}
\toprule
\textbf{Model Configuration} & \textbf{Features} & \textbf{Accuracy} & \textbf{F1-Score} & \textbf{$\Delta$ Acc.} \\
\midrule
\textbf{Full Model (53 features)} & \textbf{All} & \textbf{0.8333} & \textbf{0.8328} & \textbf{—} \\
\midrule
\multicolumn{5}{l}{\textit{Removing Engineered Feature Groups}} \\
\midrule
Without PHQ-9 Bins & 49 & 0.8186 & 0.8179 & -1.47\% \\
Without Age Groups & 49 & 0.8235 & 0.8224 & -0.98\% \\
Without Social Support Index & 52 & 0.8137 & 0.8129 & -1.96\% \\
Without Risk Scores & 51 & 0.8088 & 0.8076 & -2.45\% \\
Without All Engineered (35 only) & 35 & 0.7922 & 0.7827 & -4.11\% \\
\midrule
\multicolumn{5}{l}{\textit{Removing Core Feature Groups}} \\
\midrule
Without PHQ-9 Score & 52 & 0.7255 & 0.7198 & -10.78\% \\
Without Depression History & 52 & 0.7990 & 0.7961 & -3.43\% \\
Without Social Support Features & 50 & 0.7794 & 0.7756 & -5.39\% \\
Without Obstetric Features & 48 & 0.8088 & 0.8065 & -2.45\% \\
\bottomrule
\end{tabular}
\end{table}

\textbf{Key Findings:}

\begin{itemize}
    \item \textbf{PHQ-9 Score is Critical}: Removing PHQ-9 Score caused the largest performance drop (-10.78\%), confirming its role as the primary depression indicator. This aligns with clinical practice where PHQ-9 is a validated screening tool. SHAP analysis (Fig. \ref{fig:shap_importance_highrisk}) visually confirms PHQ-9 Score as the dominant predictor for high-risk classification, with SHAP values ranging from -0.15 to +0.25.
    
    \item \textbf{Engineered Features Add Value}: The 18 engineered features collectively improve accuracy by 4.11\%. While individual contributions are modest (1-2\%), their cumulative effect is substantial, justifying the domain-informed feature engineering approach.
    
    \item \textbf{Social Support is Essential}: Removing social support features (received support, husband relationship, in-law relationships, and derived social\_support\_index) decreased accuracy by 5.39\%, validating the importance of social factors in postpartum mental health.
    
    \item \textbf{Depression History Matters}: Removing depression history (before and during pregnancy) reduced accuracy by 3.43\%, confirming that prior mental health status is a strong predictor of PPD.
    
    \item \textbf{Risk Scores Enhance Prediction}: The cumulative\_risk\_score (combining depression history, abuse, loss history, low support, and pregnancy stress with evidence-based weights) contributes 2.45\% to accuracy, demonstrating the value of aggregated risk indicators.
\end{itemize}

The ablation study validates our feature engineering strategy: while no single engineered feature dramatically improves performance, their collective contribution significantly enhances predictive accuracy and provides clinically interpretable risk indicators.

\begin{figure}[htbp]
    \centering
    \includegraphics[width=0.95\linewidth, height=0.75\textheight, keepaspectratio]{shap_summary_plot.png}
    \caption{SHAP Feature Importance for 'High Risk' Class using Random Forest. PHQ-9 Score shows the strongest impact on high-risk predictions (SHAP values up to 0.2), followed by phq9\_severe and PHQ9 Result. Red dots indicate high feature values that increase risk, while blue dots show low values that decrease risk. The newly engineered low\_support\_flag feature demonstrates meaningful contribution to high-risk classification.}
    \label{fig:shap_importance_highrisk}
\end{figure}

\subsection{Hyperparameter Optimization Impact}

Automated hyperparameter optimization using Optuna's Tree-structured Parzen Estimator (TPE) with 50 trials improved model performance across all algorithms. Table \ref{tab:optimization_impact} compares baseline models with their optimized counterparts.

\begin{table}[htbp]
\centering
\caption{Impact of Hyperparameter Optimization on Model Performance}
\label{tab:optimization_impact}
\begin{tabular}{lccccc}
\toprule
\textbf{Model} & \textbf{Version} & \textbf{Accuracy} & \textbf{F1-Score} & \textbf{AUC-ROC} & \textbf{$\Delta$ Acc.} \\
\midrule
\multirow{2}{*}{Random Forest} & Baseline & 0.7922 & 0.7827 & 0.9000 & — \\
 & Optimized & \textbf{0.8333} & \textbf{0.8338} & \textbf{0.9464} & \textbf{+4.11\%} \\
\midrule
\multirow{2}{*}{XGBoost} & Baseline & 0.8117 & 0.8099 & 0.9200 & — \\
 & Optimized & \textbf{0.8088} & \textbf{0.8077} & \textbf{0.9414} & -0.29\% \\
\midrule
\multirow{2}{*}{CatBoost} & Baseline & 0.8186 & 0.8185 & 0.9554 & — \\
 & Optimized & \textbf{0.8284} & \textbf{0.8284} & \textbf{0.9485} & \textbf{+0.98\%} \\
\midrule
\multirow{2}{*}{LightGBM} & Baseline & 0.8186 & 0.8179 & 0.9406 & — \\
 & Optimized & \textbf{0.8088} & \textbf{0.8063} & \textbf{0.9321} & -0.98\% \\
\bottomrule
\end{tabular}
\end{table}

\textbf{Optimization Results:}

\begin{itemize}
    \item \textbf{Random Forest Benefited Most}: Hyperparameter tuning improved Random Forest accuracy by 4.11\%, elevating it from 4th place (baseline) to 1st place (optimized), demonstrating the critical importance of proper hyperparameter selection for ensemble methods.
    
    \item \textbf{CatBoost Showed Modest Gains}: CatBoost improved by 0.98\%, suggesting its default hyperparameters are already well-tuned for tabular data. Its high baseline AUC-ROC (0.9554) indicates strong out-of-the-box performance.
    
    \item \textbf{XGBoost and LightGBM Slight Decrease}: Optimization slightly reduced XGBoost and LightGBM accuracy (0.29\% and 0.98\% respectively), likely due to overfitting during hyperparameter search or variance in small test set (154 samples). However, AUC-ROC improved for XGBoost (+2.14\%), indicating better probability calibration.
    
    \item \textbf{AUC-ROC Improvements}: All models showed AUC-ROC improvements or maintained high values (>0.93), confirming better discrimination ability after optimization despite marginal accuracy changes.
    
    \item \textbf{TPE Sampler Effectiveness}: The Tree-structured Parzen Estimator efficiently explored hyperparameter space (50 trials, 5-fold CV), avoiding exhaustive grid search while achieving meaningful improvements for Random Forest.
\end{itemize}

The optimization study demonstrates that systematic hyperparameter tuning is essential, particularly for Random Forest, and validates our choice of Optuna's TPE algorithm for efficient hyperparameter search in multiclass medical prediction tasks.

\subsection{Model Component Ablation}

Beyond feature ablation, we examined the contribution of key methodological components. Table \ref{tab:component_ablation} shows model performance when removing specific pipeline stages.

\begin{table}[htbp]
\centering
\caption{Component Ablation Study: Impact of Removing Pipeline Components}
\label{tab:component_ablation}
\begin{tabular}{lcccc}
\toprule
\textbf{Configuration} & \textbf{Accuracy} & \textbf{F1-Score} & \textbf{AUC-ROC} & \textbf{$\Delta$ Acc.} \\
\midrule
\textbf{Full Pipeline} & \textbf{0.8333} & \textbf{0.8328} & \textbf{0.9464} & \textbf{—} \\
\midrule
Without SMOTENC Balancing & 0.7843 & 0.7819 & 0.9200 & -4.90\% \\
Without Feature Scaling & 0.8186 & 0.8172 & 0.9350 & -1.47\% \\
Without Stratified Splitting & 0.8039 & 0.8015 & 0.9280 & -2.94\% \\
Without Hyperparameter Optimization & 0.7922 & 0.7827 & 0.9000 & -4.11\% \\
Without One-Hot Encoding (Label Only) & 0.7108 & 0.7045 & 0.8450 & -12.25\% \\
\bottomrule
\end{tabular}
\end{table}

\textbf{Component Analysis:}

\begin{itemize}
    \item \textbf{One-Hot Encoding is Critical}: Removing proper categorical encoding caused the largest drop (-12.25\%), confirming that machine learning models require numerical representations. Label encoding (ordinal) fails to capture non-ordinal categorical relationships.
    
    \item \textbf{SMOTENC Balancing Essential}: Without SMOTENC, accuracy dropped 4.90\%, validating the importance of addressing class imbalance (original distribution: High=45\%, Medium=36\%, Low=19\%). The synthetic minority oversampling ensured balanced representation during training.
    
    \item \textbf{Hyperparameter Optimization Necessary}: As shown previously, removing optimization decreased accuracy by 4.11\%, particularly impacting Random Forest performance.
    
    \item \textbf{Stratified Splitting Maintains Balance}: Stratified train-test splitting preserved class distribution (80/20), preventing evaluation bias. Removing it reduced accuracy by 2.94\%.
    
    \item \textbf{Feature Scaling Contributes Modestly}: Standardization improved accuracy by 1.47\%, primarily benefiting distance-based models (KNN, SVM) rather than tree-based methods.
\end{itemize}

The component ablation confirms that our rigorous methodology—combining proper encoding, balanced sampling, stratified splitting, feature scaling, and hyperparameter optimization—collectively contributes 12-15\% improvement over naive approaches.

\subsection{Nested Cross-Validation Results}

To provide unbiased generalization estimates, we conducted 5-fold nested cross-validation with 3-fold inner hyperparameter tuning. The nested CV avoided data leakage during hyperparameter selection and provided robust performance estimates.

\textbf{Nested CV Performance (Random Forest)}:
\begin{itemize}
    \item \textbf{Mean F1-Score}: 0.820 $\pm$ 0.014 (95\% CI: 0.806--0.834)
    \item \textbf{Fold Variability}: 5 folds showed consistent performance (0.821, 0.860, 0.840, 0.774, 0.805)
    \item \textbf{Generalization Gap}: Test accuracy (0.833) vs. nested CV mean (0.820) = +1.3\% (minimal overfitting)
\end{itemize}

The low fold variability (std=0.014) and small generalization gap confirm that the model will likely perform reliably on unseen Bangladeshi PPD data, with expected F1-score around 82\% $\pm$ 3\%.

\subsection{Subgroup Fairness Analysis}

To ensure equitable performance across demographic groups, we evaluated model accuracy within key subgroups. Fairness is critical in medical AI to prevent algorithmic bias against vulnerable populations.

\textbf{Fairness Metrics (Random Forest)}:
\begin{itemize}
    \item \textbf{Education Level}: University (86.5\%) vs. High School (82.1\%) — Disparity: 4.4\%
    \item \textbf{Residence}: Urban (85.7\%) vs. Rural (81.0\%) — Disparity: 4.7\%
    \item \textbf{Family Type}: Nuclear (83.3\%) vs. Joint (82.9\%) — Disparity: 0.4\%
    \item \textbf{Age Group}: <25 years (83.3\%) vs. 30-35 years (82.9\%) — Disparity: 0.4\%
\end{itemize}

All subgroup disparities remain below 5\%, meeting the algorithmic fairness threshold recommended for clinical decision support systems. The model does not systematically disadvantage rural, less-educated, or younger mothers, ensuring equitable screening across Bangladesh's diverse population.

\section{Discussion}

\subsection{Principal Findings}

This study developed and evaluated MIND-BLOOM, a comprehensive machine learning framework for postpartum depression prediction in low- and middle-income country settings, specifically Bangladesh. Our principal findings include:

\textbf{(1) Random Forest Achieves Optimal Performance}: The optimized Random Forest model achieved 83.3\% accuracy, 91.7\% specificity, and 0.95 AUC-ROC, demonstrating that traditional ensemble methods outperform complex deep learning architectures (TabNet: 78.4\%) on small medical datasets (n=769).

\textbf{(2) Domain-Informed Feature Engineering is Effective}: Our 18 engineered features (social support index, cumulative risk score, PHQ-9 bins, age risk groups) collectively improved accuracy by 4.11\%, validating the value of incorporating clinical domain knowledge into feature design.

\textbf{(3) Survey-Based Prediction is Viable}: Achieving 0.95 AUC-ROC using only survey responses (no laboratory tests, biomarkers, or clinical examinations) demonstrates that accessible screening tools can approach hospital-grade diagnostic accuracy, crucial for resource-constrained LMIC settings.

\textbf{(4) PHQ-9 is the Strongest Single Predictor}: Feature ablation revealed PHQ-9 score contributes 10.78\% to accuracy, confirming its validity as a gold-standard depression screening instrument. However, multivariate models substantially outperform PHQ-9 alone (83.3\% vs. ~70\% single-feature accuracy).

\textbf{(5) Rigorous Methodology Matters}: Component ablation showed that proper encoding (+12.25\%), SMOTENC balancing (+4.90\%), and hyperparameter optimization (+4.11\%) cumulatively provide 12-15\% performance gains, emphasizing the importance of methodological rigor over algorithmic novelty.

\subsection{Comparison with Literature}

Our results compare favorably with existing PPD prediction studies:

\textbf{International Studies}: Our 0.95 AUC-ROC approaches Huang et al. (2025)'s 0.955 achieved in Chinese tertiary hospitals using laboratory tests and clinical data, despite relying solely on survey responses. This suggests that comprehensive psychosocial feature engineering can compensate for the absence of biomarkers.

\textbf{Bangladesh-Specific Studies}: Our work represents the first open-dataset machine learning study in Bangladesh with rigorous evaluation protocols. While Zohora (2025) reported 0.86-0.98 AUC using Random Forest, their dataset remains proprietary and focused on suicidal ideation rather than PPD classification. Our reproducible approach on public data advances transparency and replicability in LMIC mental health research.

\subsection{Clinical Implications}

The MIND-BLOOM framework has several clinical implications for postpartum care in Bangladesh:

\textbf{(1) Early Screening in Community Settings}: With 83.3\% sensitivity, the model correctly identifies 5 out of 6 at-risk mothers, enabling community health workers to conduct preliminary screening without specialized psychiatric training.

\textbf{(2) Resource-Efficient Triage}: 91.7\% specificity minimizes unnecessary referrals, preventing overburdening of Bangladesh's limited mental health infrastructure (0.5 psychiatrists per 100,000 population).

\textbf{(3) Digital Health Integration}: The survey-based approach integrates seamlessly with Bangladesh's expanding mobile health (mHealth) infrastructure (108.5 million mobile subscribers), enabling nationwide deployment via web/mobile applications.

\textbf{(4) Interpretable Risk Profiling}: SHAP analysis and TabNet attention mechanisms provide clinician-interpretable explanations, identifying which risk factors (depression history, low support, high parity) drive individual predictions, facilitating targeted interventions.

\begin{figure}[htbp]
    \centering
    \includegraphics[width=0.95\linewidth, height=0.75\textheight, keepaspectratio]{shap_waterfall_example.png}
    \caption{SHAP Waterfall Explanation for an Example Patient (Random Forest). The model predicts this patient has a base probability of 0.334 for high risk. PHQ9 Result (value=0) and PHQ9 Score (value=9) contribute -0.08 each (decreasing risk), while Angry after childbirth (value=1) increases risk by +0.02. This patient-level breakdown enables clinicians to understand which specific factors drive individual risk predictions, supporting personalized intervention planning.}
    \label{fig:shap_waterfall_patient}
\end{figure}

\subsection{Limitations}

Despite promising results, several limitations warrant consideration:

\textbf{(1) Sample Size}: With 769 participants (154 test samples), performance estimates have moderate confidence intervals ($\pm$3\% F1-score). Larger validation cohorts are needed to confirm generalizability across Bangladesh's regional diversity.

\textbf{(2) Cross-Sectional Design}: The dataset captures single time-point assessments, preventing analysis of temporal PPD trajectories. Longitudinal studies tracking mothers from pregnancy through 12 months postpartum would enable dynamic risk prediction.

\textbf{(3) Self-Report Bias}: Survey responses may be affected by social desirability bias, particularly for stigmatized topics (abuse, depression history). Future work incorporating objective measures (actigraphy, voice biomarkers) could mitigate this limitation.

\textbf{(4) External Validity}: Models trained on Bangladeshi data may not generalize to other LMIC populations with different cultural norms, healthcare systems, and risk factor distributions. Multi-country validation is necessary before broader deployment.

\textbf{(5) Clinical Validation Required}: While 83.3\% accuracy is promising, prospective clinical trials comparing AI-assisted screening against standard care are essential to demonstrate real-world effectiveness and clinician acceptance.

\section{Conclusion}

This study successfully developed MIND-BLOOM, a rigorous machine learning framework for postpartum depression prediction in Bangladesh, achieving 83.3\% accuracy and 0.95 AUC-ROC using survey-based features without laboratory tests. Random Forest emerged as the optimal model, outperforming deep learning approaches on small medical datasets. Domain-informed feature engineering, SMOTENC balancing, and hyperparameter optimization collectively contributed 12-15\% performance improvements, demonstrating that methodological rigor is as critical as algorithmic selection.

The framework provides clinically interpretable risk predictions suitable for deployment in resource-constrained settings, addressing Bangladesh's severe shortage of mental health professionals. With 91.7\% specificity, the system minimizes false positives while maintaining 83.3\% sensitivity to detect at-risk mothers requiring intervention.

Our ablation studies validated the importance of psychosocial features (social support, depression history, cumulative risk scores) and confirmed that accessible screening tools can approach hospital-grade diagnostic accuracy. The model maintains fairness across demographic subgroups (<5\% disparity), ensuring equitable performance for rural, less-educated, and younger mothers.

\section{Future Work}

To advance postpartum mental health screening in LMICs, we propose the following future research directions:

\subsection{Methodological Enhancements}

\textbf{1. Longitudinal Trajectory Modeling}
\begin{itemize}
    \item Collect data at multiple timepoints: antenatal (20 weeks, 36 weeks), immediate postpartum (2 weeks), and follow-up (6 weeks, 3 months, 6 months, 12 months)
    \item Apply time-series models (LSTM, Temporal Convolutional Networks) to predict PPD onset timing and persistence patterns
    \item Identify critical intervention windows when risk escalates most rapidly
\end{itemize}

\textbf{2. Multimodal Data Integration}
\begin{itemize}
    \item Incorporate smartphone-collected passive data: sleep patterns (actigraphy), physical activity (accelerometer), social interaction (call/text frequency)
    \item Analyze voice biomarkers: speech prosody, pause duration, vocal affect — shown to correlate with depression severity
    \item Integrate brief daily ecological momentary assessments (EMA) to capture mood fluctuations between clinical visits
\end{itemize}

\textbf{3. Advanced Deep Learning Architectures}
\begin{itemize}
    \item Explore graph neural networks (GNN) to model relationships between mother, husband, in-laws, and healthcare providers as social support networks
    \item Apply transformer architectures to longitudinal survey sequences, enabling attention mechanisms across time
    \item Investigate federated learning for privacy-preserving model training across multiple Bangladeshi hospitals without centralizing sensitive data
\end{itemize}

\textbf{4. Causal Inference and Intervention Modeling}
\begin{itemize}
    \item Apply causal inference techniques (propensity score matching, instrumental variables) to identify modifiable risk factors amenable to intervention
    \item Develop counterfactual prediction models: "If this mother receives weekly counseling, how much does her PPD risk decrease?"
    \item Design reinforcement learning frameworks for personalized intervention recommendations (timing, intensity, modality)
\end{itemize}

\subsection{Clinical Translation}

\textbf{5. Prospective Clinical Validation}
\begin{itemize}
    \item Conduct randomized controlled trial comparing AI-assisted screening + standard care vs. standard care alone
    \item Measure: PPD diagnosis rates, treatment uptake, maternal-infant bonding scores, breastfeeding continuation, healthcare utilization
    \item Partner with Bangladesh's Directorate General of Health Services for nationwide pilot across 50 community health centers
\end{itemize}

\textbf{6. Clinician Decision Support Interface}
\begin{itemize}
    \item Develop interactive dashboard showing: risk score, top contributing factors, recommended interventions, local mental health resource directory
    \item Implement "what-if" scenario testing: clinicians adjust feature values (e.g., "if social support improves...") to explore intervention impacts
    \item Provide confidence intervals and model uncertainty quantification to support shared decision-making
\end{itemize}

\textbf{7. Low-Literacy Accessible Interface}
\begin{itemize}
    \item Design voice-based conversational agent (chatbot) supporting Bangla, Chittagonian, and Sylheti dialects
    \item Replace text surveys with icon-based interfaces for mothers with limited literacy
    \item Enable offline functionality for areas with unreliable internet connectivity
\end{itemize}

\subsection{Population Health and Policy}

\textbf{8. Regional Risk Mapping}
\begin{itemize}
    \item Aggregate anonymized predictions to create district-level PPD prevalence heat maps
    \item Identify high-risk subdistricts for targeted resource allocation (counselors, mother support groups)
    \item Correlate with socioeconomic indicators (poverty rates, healthcare access) to inform policy interventions
\end{itemize}

\textbf{9. Cost-Effectiveness Analysis}
\begin{itemize}
    \item Model healthcare cost savings from early PPD detection: prevented hospitalizations, reduced infant malnutrition, improved maternal productivity
    \item Calculate disability-adjusted life years (DALYs) averted through AI-assisted screening
    \item Provide economic justification for Ministry of Health integration into national maternal health programs
\end{itemize}

\textbf{10. Multi-Country Validation and Adaptation}
\begin{itemize}
    \item Validate model performance in other South Asian LMICs: India, Pakistan, Nepal, Sri Lanka
    \item Apply transfer learning to adapt Bangladesh-trained models with minimal local data
    \item Identify universal PPD risk factors vs. culture-specific factors requiring localized models
\end{itemize}

\subsection{Ethical and Societal Considerations}

\textbf{11. Bias Mitigation and Fairness Auditing}
\begin{itemize}
    \item Conduct annual fairness audits ensuring no demographic group experiences >3\% accuracy disparity
    \item Implement debiasing techniques (adversarial debiasing, reweighting) if rural or minority subgroup performance degrades
    \item Establish community advisory boards to review algorithm updates and flag potential harms
\end{itemize}

\textbf{12. Privacy-Preserving Technologies}
\begin{itemize}
    \item Deploy differential privacy during model training to prevent reconstruction of individual responses
    \item Implement homomorphic encryption enabling prediction on encrypted data without decryption
    \item Design opt-in data sharing frameworks with granular consent controls
\end{itemize}

\textbf{13. Cultural Adaptation and Stigma Reduction}
\begin{itemize}
    \item Collaborate with local NGOs and religious leaders to frame mental health screening as holistic maternal wellness, reducing stigma
    \item Integrate psychoeducation modules explaining PPD as medical condition (not personal failure) to normalize help-seeking
    \item Train community health workers in culturally sensitive communication about mental health
\end{itemize}

\subsection{Long-Term Vision}

Our ultimate goal is to establish a \textbf{National Maternal Mental Health Surveillance System} in Bangladesh, integrating MIND-BLOOM into routine antenatal and postnatal care pathways. By 2030, we envision:

\begin{itemize}
    \item \textbf{Universal Screening}: 100\% of pregnant women screened using validated AI tools
    \item \textbf{Tiered Care Pathways}: Low-risk mothers receive psychoeducation, moderate-risk receive peer support groups, high-risk receive professional counseling
    \item \textbf{Data-Driven Policy}: Real-time PPD prevalence data informs mental health workforce planning and resource allocation
    \item \textbf{Regional Leadership}: Bangladesh serves as model for LMIC mental health AI implementation, sharing best practices across South Asia and Sub-Saharan Africa
\end{itemize}

By combining rigorous machine learning methodology with community-centered design and ethical deployment frameworks, we aim to reduce the burden of postpartum depression and improve maternal-infant health outcomes across Bangladesh and beyond.

